% from ega4-16-fr.tex
\oldpage[IV-4]{15}
d'anneaux $\mathcal{O}_X$-augmentés, et par suite aussi de leur limite projective
$\mathscr{P}_{X/S}^\infty$. Cet automorphisme permute les deux structures de
$\mathcal{O}_X$-Algèbre sur les $\mathscr{P}_{X/S}^n$ et sur
$\mathscr{P}_{X/S}^\infty$.
\end{proposition}

\begin{env}[16.3.5]
\label{IV.16.3.5}
Dans ce qui suit, les deux structures de $\mathcal{O}_X$-Algèbre définies sur les
$\mathscr{P}_{X/S}^n$ et sur $\mathscr{P}_{X/S}^\infty$ joueront des rôles très
différents: nous conviendrons désormais, sauf mention expresse du contraire,
que lorsque $\mathscr{P}_{X/S}^n$ ou $\mathscr{P}_{X/S}^\infty$ sera considéré
comme une $\mathcal{O}_X$-Algèbre, c'est de la structure d'Algèbre définie par
$p_1$ qu'il s'agira.
\end{env}

Pour tout ouvert $U$ de $X$ et toute section $t\in\Gamma(U,\mathcal{O}_X)$, on
notera simplement $t.1$ ou même $t$ l'image de $t$ par l'homomorphisme
structural $\Gamma(U,\mathcal{O}_X)\to\Gamma(U,\mathscr{P}_{X/S}^n)$ (resp.
$\Gamma(U,\mathcal{O}_X)\to\Gamma(U,\mathscr{P}_{X/S}^\infty)$)
(c'est-à-dire l'homomorphisme correspondant à $p_1$).

\begin{definition}[16.3.6]
\EGAReaderTarget{ega4.c04.def.16-3-6}{IV.16.3.6}
On désigne par $d_f^n$, ou $d_{X/S}^n$ (resp. $d_f^\infty$, ou
$d_{X/S}^\infty$), ou simplement $d^n$ (resp. $d^\infty$), l'homomorphisme de
faisceaux d'anneaux $\mathcal{O}_X\to\mathscr{P}_f^n=\mathscr{P}_{X/S}^n$
(resp. $\mathcal{O}_X\to\mathscr{P}_f^\infty=\mathscr{P}_{X/S}^\infty$)
déduit de $p_2$. Pour tout ouvert $U$ de $X$, et tout
$t\in\Gamma(U,\mathcal{O}_X)$, $d^nt$ (resp. $d^\infty t$) est appelé la partie
principale d'ordre $n$ (resp. partie principale d'ordre infini) de $t$. On pose
$dt=d^1t-t$, et on dit que $dt$ est la différentielle de $t$ (élément de
$\Gamma(U,\Omega_{X/S}^1)$, aussi noté $d_{X/S}(t)$).
\end{definition}

Il résulte aussitôt de cette définition que l'on a
\[
  d(t_1t_2)=t_1dt_2+t_2dt_1
  \tag{16.3.6.1}
  \label{IV.16.3.6.1}
\]
quels que soient $t_1,t_2$ dans $\Gamma(U,\mathcal{O}_X)$, autrement dit, $d$
est une dérivation de l'anneau $\Gamma(U,\mathcal{O}_X)$ dans le
$\Gamma(U,\mathcal{O}_X)$-module $\Gamma(U,\Omega_{X/S}^1)$.

Dans toutes les notations introduites dans (16.3.1) et (16.3.6), on remplacera
parfois $S$ par $A$ lorsque $S=\operatorname{Spec}(A)$.

\begin{env}[16.3.7]
\label{IV.16.3.7}
Supposons en particulier que $S=\operatorname{Spec}(A)$ et
$X=\operatorname{Spec}(B)$ soient des schémas affines, $B$ étant donc une
$A$-algèbre. Alors $\Delta_f$ correspond à l'homomorphisme canonique surjectif
$\pi:B\otimes_A B\to B$ tel que $\pi(b\otimes b')=bb'$, de noyau
$\mathfrak{I}=\mathfrak{I}_{B/A}$ (0, 20.4.1); $\mathscr{P}_f^n$ est le faisceau
structural du préschéma $\operatorname{Spec}(P_{B/A}^n)$, où
\[
  P_{B/A}^n=(B\otimes_A B)/\mathfrak{I}^{n+1};
\]
$\mathscr{G}r_\bullet(\mathscr{P}_f)$ est le $\mathcal{O}_X$-Module
quasi-cohérent correspondant au $B$-module gradué
\[
  \operatorname{gr}_{\mathfrak{I}}^\bullet(B\otimes_A B)
  =\bigoplus_{n\geq 0}(\mathfrak{I}^n/\mathfrak{I}^{n+1});
\]
en particulier $\Omega_f^1=\Omega_{X/S}^1$ est le $\mathcal{O}_X$-Module
quasi-cohérent correspondant au $B$-module des 1-différentielles de $B$ par
rapport à $A$, $\Omega_{B/A}^1$ (0, 20.4.3). Les morphismes projections
$p_1:X\times_S X\to X$, $p_2:X\times_S X\to X$ correspondent aux deux
homomorphismes d'anneaux $j_1:B\to B\otimes_A B$, $j_2:B\to B\otimes_A B$
tels que $j_1(b)=b\otimes 1$, $j_2(b)=1\otimes b$, de sorte que (par la
convention de (16.3.5)), $P_{B/A}^n$ est toujours considérée comme une
$B$-algèbre pour l'homomorphisme composé
$B\xrightarrow{j_1}B\otimes_A B\to P_{B/A}^n$; l'homomorphisme d'anneaux
$B\xrightarrow{j_2}B\otimes_A B\to P_{B/A}^n$ se note $d_{B/A}^n$ et correspond
à $d_{X/S}^n$ opérant sur $\Gamma(X,\mathcal{O}_X)$; pour tout $t\in B$, $dt$
est égal à $d_{B/A}t$, défini dans (0, 20.4.6) (cf. $\mathbf{Err}_{\mathrm{IV}}$,
11).

Si $\pi_n:B\otimes_A B\to P_{B/A}^n$ est l'homomorphisme canonique, on a donc,
en vertu des définitions précédentes
\[
  \pi_n(b\otimes b')=b.\pi_n(1\otimes b')=b.d_{B/A}^n(b')
  \quad\text{pour }b\in B,\ b'\in B.
  \tag{16.3.7.1}
  \label{IV.16.3.7.1}
\]
\end{env}

